

A closer look at the gadget constructed in
\Cref{fig:exponential-gadget} reveals a number of important structural
properties:
%
\begin{enumerate}[(i)]
\item all operations performed on the first counter are either the
  identity map $\wID$ or increments $\wPP$;
\item all operations performed on the second counter are affine
  updates $\wTT$ and $\wTTPO$;
\item once the first counter gets incremented on a run, it gets
  incremented at every subsequent transition; and
\item only counter configurations in which the value of the first
  counter equals the value of the second counter are accepted.
\end{enumerate}
%
Those properties are crucial to obtain decidability of (generalised)
existential Sem{\"e}nov arithmetic. 
%
\begin{definition}
  An \laVASS is \emph{restricted} if it has an even number of $2d$
  counters called $x_i,y_i$, $1\le i\le d$, such that every counter
  pair $(x_i,y_i)$ adheres to the above Properties~(i)--(iv), and the
  set of final counter values is defined by
  $\Phi \equiv \bigwedge_{1\le i\le d} x_i=y_i$.
\end{definition}
%
For convenience, when referring to the counters in a pair, we
subsequently refer to the first counter as its \emph{$x$-counter} and to
the second counter as its \emph{$y$-counter}. We will usually write $m$
for the value of the $x$-counter and $n$ for the value of the
$y$-counter. The following is immediate from the construction in
\Cref{prop:union-intersetion-closure}:
%
\begin{proposition}
The languages of restricted \laVASS are closed under union,
intersection, projection and inverse projections.
\end{proposition}
%
Finally, subsequently, for technical convenience, we assume that for a
restricted \laVASS, we have $\abs Q\ge 2$. This is with no loss of
generality, since if $\abs Q=1$ then deciding emptiness is trivial
(the only control state is accepting if and only if the restricted
\laVASS has non-empty language).

The next section will be devoted to the proof of the main result of
this paper on restricted \laVASS:
%
\begin{proposition}\label{prop:main}
  Language emptiness of a restricted \laVASS $V$ with $2d$ counters is
  decidable in $\mathrm{NSPACE}( \abs V \cdot 2^{O(d)})$.
\end{proposition}
%

Let us close this section with arguing how \Cref{thm:main} follows
from \Cref{prop:main}. Given a formulas $\Phi$ of generalised
Sem{\"e}nov arithmetic, we can in space $2^{O(|\Phi|)}$ construct the
disjunctive normal form of $\Phi$. Every disjunct can be assumed
to be of the form
%
\[
A\cdot \vec x = \vec b \wedge \bigwedge_{i\in I} x_i = 2^{y_i} \wedge
\bigwedge_{j\in J} R_j(\vec x),
\]
%
where the $R_j$ are predicates over regular languages. By
\Cref{lem:conjunctive-lavass}, there is a restricted \laVASS for $\Phi$ of
dimension $2 \abs I$ with a number of states bounded by $(\norm
A_{1,\infty} + \norm{\vec b}_\infty + 2)^{O(m+\abs
I)}=2^{p(\abs \Phi)}$ for some polynomial $p$ and whose language
represents the set of solutions to $A\cdot
\vec x = \vec b \wedge \bigwedge_{i\in I} x_i = 2^{y_i}$. Intersecting
with the DFA for the $R_j$ results in a restricted \laVASS $V$ with
$2\abs I=O(\abs \Phi)$ counters such that $\abs V = 2^{p(\abs \Phi)}$
for some polynomial $p$. By \Cref{prop:main}, it follows that
emptiness of $V$ is decidable in non-deterministic space exponential in
$p(\abs \Phi)$. We conclude the argument by recalling that NEXPSPACE=EXPSPACE by Savitch's theorem.
